\section{DISCUSSION AND CONCLUSION}
\label{sec:conclusion}

We presented HE-IFD, a one-shot federated knowledge distillation protocol whose central property is that every client contribution and every server-side intermediate is, by construction, a homomorphic ciphertext under collectively held keys. The server's view of training is computationally indistinguishable from random under the IND-CPA security of multiparty CKKS, and release of the trained student requires the joint participation of all clients via collective key-switching. This is a fundamentally different privacy guarantee from the statistical one offered by differential privacy: it does not trade utility for privacy, and it does not require any assumption about the composition of noise across rounds.

%Our central technical contribution is a two-phase improved training protocol. Phase 2a trains each student block independently on teacher-produced intermediate representations, providing a structured initialisation. Phase 2b refines the fully assembled student on student-produced activations, closing the train-inference distribution gap that causes block-by-block models to collapse at test time. Combined with TrainableBridge layers for per-channel alignment, client-side collaborative normalisation, and magnitude-regularised loss, this protocol achieves stable convergence where earlier approaches diverge.



Our experiments across three architectures (ResNet-18, SimpleCNN, and ViT-B/32) confirm that the framework generalises beyond a single model family.
Compared against differential privacy baselines, HE-IFD achieves competitive accuracy with a strictly stronger privacy guarantee: at $N$=4, $\alpha=0.1$, HE-IFD reaches 57.4\% while the best DP-SGD configuration (epsilon=10, weakest privacy) achieves 64.2\%, and tighter DP settings ($\varepsilon$ $\leq$ 1) collapse accuracy below 33\%.

Our privacy analysis demonstrates that standard membership inference attacks (loss-threshold and confidence-based) achieve near-random AUC (${\sim}$0.5) against the distilled student. As reviewed in Section~\ref{sec:bg_server_inference}, stronger attacks from the literature, e.g., image reconstruction from logits, feature inversion, and advanced likelihood-ratio attacks, all require access to plaintext knowledge transfers, which HE eliminates entirely.

Future work along the present threat model includes reducing the one-shot upload cost (${\sim}$115\,GB per client at $N{=}4$, $\sim 460$ GB total across all clients) through ciphertext compression and feature selection, and investigating bootstrapping-free training schedules for deeper HE-compatible architectures.

A second class of future work concerns adversaries that fall outside the present threat model. The semi-honest assumption on clients bounds the contribution to passive collusion; an actively malicious client can craft adversarial logit ciphertexts (encrypted-feature poisoning) or submit ciphertexts engineered to bias the aggregate $\widetilde{Y}$ (model poisoning under encryption), neither of which the binding invariant of Section~\ref{sec:bg_server_inference} addresses. The classical plaintext defences against these threats---coordinate-wise median aggregation, Krum, trimmed mean---do not directly port to the homomorphic setting, since their rank-based statistics are not realisable as low-depth arithmetic circuits on ciphertexts. A robust aggregation primitive compatible with HE is therefore an open problem distinct from the present work's contribution. The natural extension toward verifiable correctness---independent of robustness---is verifiable HE: a server proves in zero knowledge that the computation it performed on ciphertexts was the prescribed one, as surveyed by~\cite{viand2023verifiable} and concretised for CKKS by~\cite{atapoor2024vfhe}. We do not claim defence against any of these adversaries; we acknowledge the threat surface explicitly so that the contribution of HE-IFD is correctly bounded to what is established here: privacy of the client data against a semi-honest server colluding with up to $N{-}1$ clients. %Our results at $\alpha{=}10.0$ (near-homogeneous) confirm a performance ceiling for knowledge-diversity gains: at $N{=}4$, the student falls $-8.5$ pp below the best individual teacher; at $N{=}16$ it reaches $-2.9$ pp below the best but $+2.4$ pp above the mean. When all clients hold near-IID data, HE-IFD remains useful for privacy---encrypting the upload eliminates MIA attack surface regardless of data distribution---but the accuracy advantage over individual teachers that motivates participation in the heterogeneous setting is substantially reduced.

\section*{Acknowledgments}
We acknowledge the support of TÜBİTAK (the Scientific and Technological Research Council of Türkiye) projects 124E091 and 124N941. The authors used Claude Sonnet 4.6 and Claude Opus 4.6 to enhance the manuscript by shortening sentences, correcting typos, and improving grammar.